\documentclass[10pt,a4paper]{amsart}
\usepackage{amsmath,amssymb,amsthm,graphicx}
\usepackage{verbatim}
\newenvironment{reference}{\comment}{\endcomment}
%\newenvironment{reference}{}{}
\newenvironment{slogan}{\comment}{\endcomment}
\newenvironment{history}{\comment}{\endcomment}

% For commutative diagrams we use Xy-pic
\usepackage[all]{xy}

% We use 2cell for 2-commutative diagrams.
\xyoption{2cell}
% Theorem environments.
%
\theoremstyle{plain}
\newtheorem{theorem}[subsection]{Theorem}
\newtheorem{proposition}[subsection]{Proposition}
\newtheorem{lemma}[subsection]{Lemma}

\theoremstyle{definition}
\newtheorem{definition}[subsection]{Definition}
\newtheorem{example}[subsection]{Example}
\newtheorem{exercise}[subsection]{Exercise}
\newtheorem{situation}[subsection]{Situation}

\theoremstyle{remark}
\newtheorem{remark}[subsection]{Remark}
\newtheorem{remarks}[subsection]{Remarks}

\numberwithin{equation}{subsection}

% Macros
%
\def\lim{\mathop{\mathrm{lim}}\nolimits}
\def\colim{\mathop{\mathrm{colim}}\nolimits}
\def\Spec{\mathop{\mathrm{Spec}}}
\def\Hom{\mathop{\mathrm{Hom}}\nolimits}
\def\Ext{\mathop{\mathrm{Ext}}\nolimits}
\def\SheafHom{\mathop{\mathcal{H}\!\mathit{om}}\nolimits}
\def\SheafExt{\mathop{\mathcal{E}\!\mathit{xt}}\nolimits}
\def\Sch{\mathit{Sch}}
\def\Mor{\mathop{\mathrm{Mor}}\nolimits}
\def\Ob{\mathop{\mathrm{Ob}}\nolimits}
\def\Sh{\mathop{\mathit{Sh}}\nolimits}
\def\NL{\mathop{N\!L}\nolimits}
\def\CH{\mathop{\mathrm{CH}}\nolimits}
\def\proetale{{pro\text{-}\acute{e}tale}}
\def\etale{{\acute{e}tale}}
\def\QCoh{\mathit{QCoh}}
\def\Ker{\mathop{\mathrm{Ker}}}
\def\Im{\mathop{\mathrm{Im}}}
\def\Coker{\mathop{\mathrm{Coker}}}
\def\Coim{\mathop{\mathrm{Coim}}}

% Boxtimes
%
\DeclareMathSymbol{\boxtimes}{\mathbin}{AMSa}{"02}

%
% Macros for moduli stacks/spaces
%
\def\QCohstack{\mathcal{QC}\!\mathit{oh}}
\def\Cohstack{\mathcal{C}\!\mathit{oh}}
\def\Spacesstack{\mathcal{S}\!\mathit{paces}}
\def\Quotfunctor{\mathrm{Quot}}
\def\Hilbfunctor{\mathrm{Hilb}}
\def\Curvesstack{\mathcal{C}\!\mathit{urves}}
\def\Polarizedstack{\mathcal{P}\!\mathit{olarized}}
\def\Complexesstack{\mathcal{C}\!\mathit{omplexes}}
% \Pic is the operator that assigns to X its picard group, usage \Pic(X)
% \Picardstack_{X/B} denotes the Picard stack of X over B
% \Picardfunctor_{X/B} denotes the Picard functor of X over B
\def\Pic{\mathop{\mathrm{Pic}}\nolimits}
\def\Picardstack{\mathcal{P}\!\mathit{ic}}
\def\Picardfunctor{\mathrm{Pic}}
\def\Deformationcategory{\mathcal{D}\!\mathit{ef}}
\usepackage{cleveref}
\crefname{originallemma}{original lemma}{original lemmas}
\Crefname{originallemma}{Original lemma}{Original lemmas}
\makeatletter
\newlabel{algebra-lemma-qc-open}{{00F6}{source}{}{sourcepointer-1}{}}
\newlabel{algebra-lemma-qc-open@cref}{{[originallemma][0][]00F6}{[1][0][]source}{}{}{}}
\makeatother
\crefname{originallemma}{original lemma}{original lemmas}
\Crefname{originallemma}{Original lemma}{Original lemmas}
\makeatletter
\newlabel{algebra-lemma-support-element}{{07Z5}{source}{}{sourcepointer-2}{}}
\newlabel{algebra-lemma-support-element@cref}{{[originallemma][0][]07Z5}{[1][0][]source}{}{}{}}
\makeatother
\crefname{originaldefinition}{original definition}{original definitions}
\Crefname{originaldefinition}{Original definition}{Original definitions}
\makeatletter
\newlabel{modules-definition-support}{{01AT}{source}{}{sourcepointer-3}{}}
\newlabel{modules-definition-support@cref}{{[originaldefinition][0][]01AT}{[1][0][]source}{}{}{}}
\makeatother
\crefname{originallemma}{original lemma}{original lemmas}
\Crefname{originallemma}{Original lemma}{Original lemmas}
\makeatletter
\newlabel{schemes-lemma-push-forward-quasi-coherent}{{01LC}{source}{}{sourcepointer-4}{}}
\newlabel{schemes-lemma-push-forward-quasi-coherent@cref}{{[originallemma][0][]01LC}{[1][0][]source}{}{}{}}
\makeatother
\crefname{originalsection}{original section}{original sections}
\Crefname{originalsection}{Original section}{Original sections}
\makeatletter
\newlabel{schemes-section-quasi-coherent}{{01LA}{source}{}{sourcepointer-5}{}}
\newlabel{schemes-section-quasi-coherent@cref}{{[originalsection][0][]01LA}{[1][0][]source}{}{}{}}
\makeatother
\title{Sections supported on a closed subset with retrocompact complement form a quasi-coherent sheaf}
\author{The Stacks Project Authors}
\date{stacks-project:properties.tex:local-snapshot-20261012; GNU FDL 1.2 or later}
\begin{document}
\maketitle
\noindent\textit{Editorial note.} All delivered mathematical text is the author's own, copied from the pinned Stacks Project chapter file properties.tex; no statement or proof step was rewritten, reordered or removed.\par\medskip
\noindent\textit{Editorial note.} The counted proof is the author's complete argument: the finitely generated ideals provided by retrocompactness, the translation of the support condition into annihilation by a power of the ideal through the support lemma, and the identification of the sheaf as a kernel of a map to a quasi-coherent pushforward.\par\medskip
\noindent\textit{Editorial note.} The label lemma-sections-supported-on-closed-subset names the counted statement, and the proof environment immediately following it in the source is delivered verbatim as its proof.\par\medskip
\noindent\textit{Editorial note.} The counted statement is carried by the author's own lemma environment in the chapter source.\par\medskip
\noindent\textit{Editorial note.} Pre-open gate: python operations/agent-authoring/tools/precard\_gate.py --source properties.tex --label lemma-sections-supported-on-closed-subset --start 4051 --end 4108 --local-id 28038 returned RESULT=OK (L28038 free, occupancy FREE, public interval clear). Candidate taken from sources/stacks-properties/candidates-halfA.json.\par\medskip
\noindent\textit{Editorial note.} The statement and proof are copied from the pinned chapter file; the cited inputs are kept as external pointers with their canonical Stacks tags.\par\medskip
\noindent\textit{Allowed standard prerequisites (editorial).}\begin{itemize}
\item Finiteness of the ideal defining a closed subset in an affine open. Algebra, the lemma that a closed subset of an affine scheme with retrocompact complement is defined by a finitely generated ideal. Supplies the finitely generated ideal I with Z cap U = V(I).
\item Description of the support of a section and annihilation by the ideal. Modules, the definition of the support of a section, and Algebra, the lemma relating the vanishing of a localization to the containment of the support in a closed subset. Supplies the equivalence between the support of x being contained in V(I) and the annihilation of x by some power of I.
\item Quasi-coherence of the pushforward along a quasi-compact open immersion. Schemes, the lemma that the pushforward along a quasi-compact and separated morphism of a quasi-coherent module is quasi-coherent. Supplies the quasi-coherence of the pushforward of the restriction to the complement, so that the kernel of the displayed map is quasi-coherent.
\end{itemize}\medskip
\ifdefined\hypertarget\hypertarget{sourcepointer-1}{}\fi
\noindent\textit{Source reference.} Original lemma 00F6: sources/stacks-chow/tags.txt, source lines 560--560. Prerequisite pointer invoked at proof line(s) 4075; it names a standard fact of a sibling chapter that the counted argument uses, and every step of the counted proof built on it is retained verbatim in the delivered body.\par\medskip
\ifdefined\hypertarget\hypertarget{sourcepointer-2}{}\fi
\noindent\textit{Source reference.} Original lemma 07Z5: sources/stacks-chow/tags.txt, source lines 10135--10135. Prerequisite pointer invoked at proof line(s) 4085; it names a standard fact of a sibling chapter that the counted argument uses, and every step of the counted proof built on it is retained verbatim in the delivered body.\par\medskip
\ifdefined\hypertarget\hypertarget{sourcepointer-3}{}\fi
\noindent\textit{Source reference.} Original definition 01AT: sources/stacks-chow/tags.txt, source lines 1702--1702. Prerequisite pointer invoked at proof line(s) 4083; it names a standard fact of a sibling chapter that the counted argument uses, and every step of the counted proof built on it is retained verbatim in the delivered body.\par\medskip
\ifdefined\hypertarget\hypertarget{sourcepointer-4}{}\fi
\noindent\textit{Source reference.} Original lemma 01LC: sources/stacks-chow/tags.txt, source lines 2081--2081. Prerequisite pointer invoked at proof line(s) 4103; it names a standard fact of a sibling chapter that the counted argument uses, and every step of the counted proof built on it is retained verbatim in the delivered body.\par\medskip
\ifdefined\hypertarget\hypertarget{sourcepointer-5}{}\fi
\noindent\textit{Source reference.} Original section 01LA: sources/stacks-chow/tags.txt, source lines 2079--2079. Prerequisite pointer invoked at proof line(s) 4107; it names a standard fact of a sibling chapter that the counted argument uses, and every step of the counted proof built on it is retained verbatim in the delivered body.\par\medskip
\begin{lemma}
\label{lemma-sections-supported-on-closed-subset}
Let $X$ be a scheme. Let $Z \subset X$ be a closed subset.
Let $\mathcal{F}$ be a quasi-coherent $\mathcal{O}_X$-module.
Consider the sheaf of $\mathcal{O}_X$-modules $\mathcal{F}'$
which associates to every open $U \subset X$
$$
\mathcal{F}'(U)
=
\{s \in \mathcal{F}(U) \mid
\text{the support of }s\text{ is contained in }Z \cap U\}
$$
If $X \setminus Z$ is a retrocompact open of $X$, then
\begin{enumerate}
\item for an affine open $U \subset X$ there exist a finitely generated
ideal $I \subset \mathcal{O}_X(U)$ such that $Z \cap U = V(I)$,
\item for $U$ and $I$ as in (1) we have
$\mathcal{F}'(U) = \{x \in \mathcal{F}(U) \mid
I^nx = 0 \text{ for some } n\}$,
\item $\mathcal{F}'$ is a quasi-coherent sheaf of $\mathcal{O}_X$-modules.
\end{enumerate}
\end{lemma}

\begin{proof}
Part (1) is Algebra, Lemma \ref{algebra-lemma-qc-open}.
Let $U = \Spec(A)$ and $I$ be as in (1).
Then $\mathcal{F}|_U$ is the quasi-coherent sheaf associated to
some $A$-module $M$. We have
$$
\mathcal{F}'(U) = \{x \in M \mid x = 0\text{ in }M_\mathfrak p
\text{ for all }\mathfrak p \not \in Z\}.
$$
by Modules, Definition \ref{modules-definition-support}. Thus
$x \in \mathcal{F}'(U)$ if and only if $V(\text{Ann}(x)) \subset V(I)$, see
Algebra, Lemma \ref{algebra-lemma-support-element}. Since $I$ is
finitely generated this is equivalent to $I^n x = 0$ for some $n$.
This proves (2).

\medskip\noindent
Proof of (3). Observe that given $U \subset X$ open there is an exact sequence
$$
0 \to \mathcal{F}'(U) \to \mathcal{F}(U) \to \mathcal{F}(U \setminus Z)
$$
If we denote $j : X \setminus Z \to X$ the inclusion morphism, then
we observe that $\mathcal{F}(U \setminus Z)$ is the sections of the
module $j_*(\mathcal{F}|_{X \setminus Z})$ over $U$. Thus we have
an exact sequence
$$
0 \to \mathcal{F}' \to \mathcal{F} \to j_*(\mathcal{F}|_{X \setminus Z})
$$
The restriction $\mathcal{F}|_{X \setminus Z}$ is quasi-coherent.
Hence $j_*(\mathcal{F}|_{X \setminus Z})$ is quasi-coherent
by Schemes, Lemma \ref{schemes-lemma-push-forward-quasi-coherent}
and our assumption that $j$ is quasi-compact (any open immersion
is separated). Hence $\mathcal{F}'$ is quasi-coherent
as a kernel of a map of quasi-coherent modules, see
Schemes, Section \ref{schemes-section-quasi-coherent}.
\end{proof}

\end{document}
